% Generated draft: requires mathematical and source review.
\documentclass[11pt]{article}
\usepackage[margin=25mm]{geometry}
\usepackage{fontspec}
\setmainfont{Latin Modern Roman}
\usepackage{amsmath,amssymb,mathtools,unicode-math}
\setmathfont{Latin Modern Math}
\usepackage{enumitem,needspace,xurl,hyperref}
\setlength{\parindent}{0pt}
\setlength{\parskip}{5pt}
\setlength{\emergencystretch}{3em}
\newcommand{\N}{\mathbb{N}}
\newcommand{\Z}{\mathbb{Z}}
\newcommand{\Q}{\mathbb{Q}}
\newcommand{\R}{\mathbb{R}}
\newcommand{\C}{\mathbb{C}}
\newcommand{\F}{\mathbb{F}}
\newcommand{\A}{\mathbb{A}}
\newcommand{\PP}{\mathbb{P}}
\newcommand{\E}{\mathbb{E}}
\newcommand{\eps}{\varepsilon}
\newcommand{\dd}{\,\mathrm{d}}
\begin{document}
\section*{UnsolvedMath ID: 30006379}
\textbf{Source problem code:} catalogue entry
\subsection{Definitions and mathematical statement}
The supplied record contains the rhetorical heading ``Why should we replace Galois cohomology?'' followed by the talk title ``Duality for condensed cohomology of the Weil group of a \(p\)-adic field'' and the name Marco Artusa.

No mathematical problem can be recovered precisely from this text. It supplies no definitions of the relevant cohomology theories, coefficient objects, hypotheses, comparison maps, duality pairing, or quantified assertion. In particular, the heading does not specify a proposed replacement for Galois cohomology or a property that such a replacement should satisfy. The talk title identifies a topic but does not state a duality theorem or conjecture.

The supplied statement status is \texttt{unrecoverable}. According to the supplied statement-verification note, the extraction is a rhetorical talk heading followed by its title and author, rather than a standalone open problem; no clean statement is promoted. Thus this entry preserves the extraction and its provenance without introducing a mathematical assertion. The registry label \texttt{open} is a supplied classification, not a verified claim that a particular conjecture remains unresolved.
\subsection{Short English statement}
This record preserves a rhetorical heading and a talk title attributed to Marco Artusa. It does not specify a mathematical question about replacing Galois cohomology or a duality statement for condensed cohomology. A precise problem statement requires further source material and human mathematical review.
\subsection{Sources}
\begin{itemize}
\item Local UnsolvedMath registry, permanent ID \(30006379\), source code \texttt{catalogue entry}, title ``Rhetorical Condensed-Cohomology Talk Heading.'' The supplied record belongs to the Oberwolfach Reports Open Problems collection and retains the original extraction for provenance.
\item Supplied report citation: \emph{Algebraic K-theory (2026)}, Oberwolfach Reports, DOI \texttt{10.4171/owr/2025/36}, pp.\ 1949--2012. Supplied source URL and literature reference labelled ``Oberwolfach report'': \url{https://doi.org/10.4171/owr/2025/36}. The supplied proposed year and workshop-year tag are 2025; their relationship to the citation's 2026 designation is not explained.
\item Marco Artusa is named in the extracted talk heading. The record supplies no separate problem proposer and no authoritative citation to a standalone mathematical problem. Its statement-verification note reports that the official report identifies this extraction as a rhetorical heading, talk title, and author; that verification is supplied metadata, not an independent source check performed here.
\item Supplied literature assessment, checked on exactly \texttt{2026-08-22}: status ``open''; the rhetorical question has no precise replacement theory in the extracted record, and no uniquely identifiable resolved claim can be assigned to it. This is supplied dated curation, not independently verified. The supplied background requests independent verification of time-sensitive status and cited claims.
\end{itemize}
\end{document}
